% 11-conclusions.tex
\chapter{Conclusions}
\label{ch:conclusions}

This report has presented Doki, a complete container ecosystem designed
for Android and Linux environments. Doki addresses the fundamental
challenge of running containers on mobile devices by using proot-based
syscall interception, filesystem detection probes, and adaptive
networking to operate within the constraints of Android's restricted
kernel.

\section{Summary of Contributions}
\label{sec:contributions-summary}

The key contributions of this work include:

\begin{enumerate}[leftmargin=*, itemsep=4pt]
  \item \textbf{Twelve-level isolation framework.}
        (\cref{sec:isolation-levels}) A runtime registry that
        automatically selects the strongest available isolation
        mechanism, from hardware-enforced microVMs to lightweight
        syscall interception via proot. The framework provides both
        security and compatibility across diverse hardware
        capabilities.

  \item \textbf{DokiLink-Lite mesh networking.}
        (\cref{sec:dokilink}) A peer-to-peer mesh networking protocol
        with TLS~1.3 encryption and optional NaCl secretbox payload
        encryption, enabling secure inter-device communication without
        centralized infrastructure.

  \item \textbf{Docker-compatible build and orchestration.}
        (\cref{ch:build}) A build system and Compose engine
        that accepts standard Dockerfile syntax and Compose YAML,
        lowering the barrier to adoption for developers familiar with
        the Docker ecosystem.

  \item \textbf{Empirical evaluation.}
        (\cref{ch:metrics}) A comprehensive performance analysis
        demonstrating that proot-based container isolation adds
        acceptable overhead (100--150\,ms startup, 15\,MB memory per
        container) for development, education, and edge computing use
        cases.
\end{enumerate}

\section{Key Metrics}
\label{sec:key-metrics}

The following table summarizes the key metrics achieved by Doki v0.9.3:

\begin{table}[htbp]
  \centering
  \caption{Summary of Doki v0.9.3 key metrics.}
  \label{tab:summary-metrics}
  \begin{tabularx}{\textwidth}{l X}
    \toprule
    \textbf{Metric} & \textbf{Value} \\
    \midrule
    CLI commands          & 244 across 9 categories \\
    Isolation levels      & 12 (microVM to native) \\
    Binary size           & 28.3\,MB total (4 statically linked binaries) \\
    Startup overhead      & 100--150\,ms vs.\ Docker on server \\
    Memory per container  & 15\,MB additional overhead \\
    Network encryption    & TLS 1.3 + optional NaCl secretbox \\
    Architectures         & ARM64, ARMv7 (Android, Linux, macOS) \\
    Test coverage         & 14 test packages, all assertions passing \\
    \bottomrule
  \end{tabularx}
\end{table}

\section{Limitations}
\label{sec:limitations}

Doki has several known limitations that should be considered:

\begin{itemize}[leftmargin=*, itemsep=4pt]
  \item \textbf{Performance overhead.} The proot-based syscall
        interception adds measurable overhead compared to
        server-class Docker. Container startup times are 2--3x slower,
        and memory usage is approximately 15\,MB higher per container.

  \item \textbf{Isolation constraints.} Without root access, Doki
        cannot provide kernel-level namespace isolation on most Android
        devices. The proot runtime relies on \texttt{ptrace()}, which
        is slower than native namespaces and does not provide full
        process isolation.

  \item \textbf{Networking limitations.} The absence of
        \texttt{iptables} on Android requires fallback to socat-based
        port forwarding, which is less efficient than kernel-level NAT.
        Port 53 is unavailable due to SELinux restrictions.

  \item \textbf{Health check execution.} Compose
        \texttt{service\_healthy} conditions are parsed but not yet
        executed by the container runtime. This limits the
        reliability of multi-service deployments.

  \item \textbf{Platform dependency.} Doki relies on
        \texttt{proot}~\cite{proot2024} as its primary isolation
        mechanism. If proot becomes unavailable or incompatible with
        future Android versions, alternative isolation backends would
        need to be developed.
\end{itemize}

\section{Future Impact}
\label{sec:future-impact}

The container ecosystem on mobile devices remains an emerging field.
Doki demonstrates that Docker-like functionality is achievable without
root access, kernel modifications, or cloud connectivity. As mobile
hardware continues to gain processing power and memory, the demand
for local containerization will grow.

Potential applications of Doki include:

\begin{itemize}[leftmargin=*, itemsep=4pt]
  \item \textbf{Mobile development environments} --- running
        compilers, linters, and test suites directly on Android
        devices, eliminating the need for cloud-based CI/CD for
        quick iterations.

  \item \textbf{Edge computing} --- deploying containerized services
        on mobile devices at the network edge, enabling low-latency
        processing for IoT and sensor data.

  \item \textbf{Education} --- teaching container concepts using
        widely available Android hardware, making cloud-native
        technologies accessible to students without server access.

  \item \textbf{Field operations} --- running isolated tools in
        environments without traditional server infrastructure,
        such as remote sensing, disaster response, and military
        operations.
\end{itemize}

Doki provides a foundation for these use cases. Future development
will focus on strengthening isolation (gVisor integration), improving
networking (WireGuard mesh, NAT traversal), and enabling ecosystem
integration (Kubernetes CRI compliance).
